\section*{Résumé}
\addcontentsline{toc}{section}{Résumé}
\markboth{Résumé}{}

% Résumé du travail en français (200--300 mots) :
% contexte, objectifs, méthode utilisée et principaux résultats.

\begin{otherlanguage}{french}
  % French text here
Les processus traditionnels de migration de données, notamment lors de la phase d'analyse et de conception, souffrent d'une forte dépendance aux tâches manuelles et d'une fragmentation des outils de gestion (tableurs, documents bureautiques décentralisés). Cette approche classique, sujette aux erreurs humaines, entraîne des délais prolongés et un manque de traçabilité limitant l'efficacité globale des projets.

L'objectif de ce projet est de concevoir et développer l'ADMP AI Studio, une plateforme intelligente visant à moderniser l'architecture système et à automatiser les pipelines ETL (Extract, Transform, Load) de migration de données au sein du groupe Accenture. Le but est de réduire drastiquement l'intervention manuelle tout en standardisant la production des règles de transformation.

La méthode adoptée repose sur le déploiement d'une architecture novatrice d'Intelligence Artificielle agentique. Un système multi-agents, orchestré via LangGraph et intégrant des modèles de langage avancés (LLMs) utilisant le raisonnement "Chain of Thought", a été mis en place. Techniquement, l'architecture transitionne vers un stockage centralisé sous PostgreSQL exploitant le format JSON. Différents agents spécialisés (Scope, Design, Validation) travaillent en synergie pour ingérer les modèles de données, effectuer des analyses d'écarts (Gap Analysis) de manière autonome, et valider la cohérence des opérations.

Les principaux résultats mettent en évidence l'automatisation réussie de la génération des livrables critiques de conception (Data Mapping, Table Mapping, SDS, TDS). La nouvelle infrastructure établit un flux de données continu, sécurisé et auditable, fiabilisant ainsi la qualité des données avant leur injection vers les systèmes cibles, et préparant le terrain pour l'automatisation des phases de transformation complexes.


\end{otherlanguage}


\vspace{1em}
\noindent\textbf{Mots-clés :} Migration de données, Intelligence Artificielle agentique, Pipeline ETL, Architecture système, Automatisation.
